\section{Small traces remove properness parameters}
\label{annex:p06-properness-parameter-invariance}

\begin{annexitem}{annex:p06.ppi.small-absoluteness}{thm:p06-ppi-small-absoluteness}{Hereditary bounds and the exact trace predicate}
\begin{proof}
The transitive closure of \(\mathbf P\) has size at most \(\kappa\),
so every member of \(P\), and every subset of \(P\), has hereditary
size at most \(\kappa\). The whole \(X\) has hereditary size at most
\(2^\kappa\). Countable subsets of \(X\) have hereditary size at most
\(\kappa\); their enumerations have that same bound. There are at
most \(|X|^{\aleph_0}\leq(2^\kappa)^{\aleph_0}=2^\kappa\) such
subsets, using Choice and \(\kappa\cdot\aleph_0=\kappa\).
Consequently their complete collection is hereditarily at most
\(2^\kappa\), as is any subcollection \(C\) of it.

For a code \(g\), the finite-input domain has size at most \(|X|\),
and every countable output has hereditary size at most \(\kappa\).
Its entire graph therefore has hereditary size at most \(2^\kappa\).
Every object just listed belongs to \(H_\chi\) under the strict bound
\(\chi>2^\kappa\). Every actually countable subset and every increasing
countable sequence from \(C\) also has its enumeration/graph in
\(H_\chi\). Thus countability, unboundedness and increasing-union
closure are tested against exactly the same objects in both structures.
The code-closure predicate has only bounded finite-input/output quantifiers.

The density predicate for \(D\subseteq P\), and the compatibility
predicate, quantify only over \(P\), \(D\), and the fixed order.
The good-trace formula likewise quantifies over \(Y\), \(\mathcal D_P\)
and \(P\); these full sets are present. Transitivity makes those
bounded assertions absolute. This proves the asserted absoluteness
without a Power Set axiom for arbitrary \(H_\chi\)-elements.
In particular it does not need the set of all possible codes, which
can be much larger than \(H_\chi\).

The internal power set of \(P\) agrees with the actual one because
every subset is present, and that power set itself is present. Thus
\(X\) is uniquely definable from the poset datum in \(H_\chi\).
For \(N\prec H_\chi\), it belongs to \(N\), and
\(\omega\subseteq N\) by \cref{thm:p06-pfm-hulls}. Hence
\(Y=N\cap X\) is countably infinite. Since \(P\subseteq X\) and
every dense \(D\subseteq P\) lies in \(\mathcal P(P)\subseteq X\),
\[
 P\cap Y=P\cap N,\quad \mathcal D_P\cap Y=\mathcal D_P\cap N,
 \quad D\cap Y=D\cap N.
\]
Substitution gives exactly the master-production formula, including
its single existential \(q\) preceding all dense-set tests.
\end{proof}
\end{annexitem}

\begin{annexitem}{annex:p06.ppi.club-codes}{thm:p06-ppi-club-codes}{Build closure clubs and code arbitrary clubs}
\begin{proof}
For a given code \(g\), start with a countably infinite set \(Y_0\)
containing any requested countable seed. Recursively add every
\(g(a)\) for finite \(a\subseteq Y_n\). Each stage is countable:
it has countably many finite inputs and countable outputs. Their union
\(Y\) is countably infinite and closed under \(g\), since each finite
input in the union occurs in one stage. This proves unboundedness.
The union of an increasing countable chain of closed sets is closed
for the same finite-input reason, proving clubhood.

Conversely, fix club \(C\). Recursively on the finite size of \(a\),
choose \(c_a\in C\) containing \(a\) and all \(c_b\) for proper
subsets \(b\subsetneq a\). At the empty input choose any \(c_\varnothing\in C\).
At each other input there are only finitely many previous countable
sets to include, so unboundedness of \(C\) supplies a witness. Choice
selects the whole set-indexed family. Put \(g(a)=c_a\).
If \(Y\) is countably infinite and closed under this \(g\), enumerate
it without repetitions and let \(a_n\) be the first \(n+1\) elements.
The \(c_{a_n}\) form an increasing sequence of members of \(C\),
all inside \(Y\), whose union is \(Y\). Closedness puts \(Y\) in
\(C\). Therefore \(C_g\subseteq C\), proving the converse fully.
\end{proof}
\end{annexitem}

\begin{annexitem}{annex:p06.ppi.trace-projection}{thm:p06-ppi-trace-projection}{Project by finite-character hulls, not a naive image}
\begin{proof}
Apply the club-code lemma on the set \(Z=H_\Theta\) to the given
large-model club \(C\), obtaining a countable-output code \(f\) with
\(C_f\subseteq C\). Let \(h(S)\) be the least \(f\)-closed subset
of \(Z\) containing \(S\), formed by the countable closure construction.
Even a finite seed has countably infinite closure, since it includes
\(f(\varnothing)=c_\varnothing\in C\).

For a finite \(a\subseteq X\), define \(g(a)=h(a)\cap X\).
The finite-character identity is
\[
 h(Y)=\bigcup\{h(a):a\in[Y]^{<\omega}\}.
\]
Indeed the right side contains \(Y\); for a finite subset of it,
choose finitely many witnessing finite seeds and take their union.
Its hull contains that subset and all of its \(f\)-outputs. Thus the
right side is closed and is exactly the least closure.

For \(Y\in C_g\), this identity gives \(h(Y)\cap X\subseteq Y\),
while the reverse inclusion is immediate. The hull is countable and
\(f\)-closed, so belongs to \(C\). Write \(M=h(Y)\); by hypothesis
it is elementary, contains the poset and has master production.
For \(p\in Y\cap P=M\cap P\), choose its \(M\)-master \(q\).
For each dense \(D\in Y=M\cap X\), one has
\(D\cap Y=D\cap M\), so that same condition satisfies every
good-trace test. Hence \(C_g\) is a club of good traces.

Although \(f\) and \(h(a)\) may involve arbitrarily large parameters
or models, \(g\) has inputs in \(X\) and countable outputs in \(X\).
The small-absoluteness lemma bounds its hereditary size by
\(2^\kappa\). This is the exact mechanism for removing the large datum.
\end{proof}
\end{annexitem}

\begin{annexitem}{annex:p06.ppi.all-model-recognition}{thm:p06-ppi-all-model-recognition}{Selecting a small witness inside the tested model}
\begin{proof}
Take a club of good traces and apply the club-code lemma to obtain
\(g:[X]^{<\omega}\to[X]^{\leq\omega}\) such that
\(C_g\subseteq\{Y:\operatorname{Good}_P(Y)\}\). Fix any regular
\(\chi>2^\kappa\). The individual code belongs to \(H_\chi\).
Let \(\Phi(\mathbf P,g)\) be the following first-order assertion:
\(g\) is a function with the displayed domain and range, and every
\(Y\in[X]^\omega\) closed under its outputs satisfies
\(\operatorname{Good}_P(Y)\). Here \(X\) is the uniquely specified
set \(\omega\cup P\cup\mathcal P(P)\), not an extra parameter.
By the small-absoluteness lemma, the actual code witnesses
\(H_\chi\models\exists g\,\Phi(\mathbf P,g)\). The assertion uses
the full actual sets of countable subsets and dense subsets, which are
present at this height; it does not quantify only over those objects
already in a particular elementary model.

Now let \(N\prec(H_\chi,\in)\) be any countable model with
\(\mathbf P\in N\). Elementarity selects a witness
\(g_N\in N\) with \(H_\chi\models\Phi(\mathbf P,g_N)\).
There is no claim that \(g_N\) is the original code. Set
\(Y=N\cap X\). This trace is countably infinite. If
\(a\in[Y]^{<\omega}\), finite-set closure of elementary models
gives \(a\in N\). The unique value \(g_N(a)\) belongs to \(N\).
It is at most countable, so the countable-element closure in
\cref{thm:p06-pfm-hulls} gives \(g_N(a)\subseteq N\).
For an empty value this inclusion is immediate; otherwise an
enumeration in \(N\), evaluated at every natural number in \(N\),
proves it. Since the value is also a subset of \(X\), it is contained
in \(Y\). Thus \(Y\in C_{g_N}\).

The assertion \(\Phi\), and its good-trace consequence for this
actual \(Y\), are absolute at the stated height. Therefore
\(\operatorname{Good}_P(Y)\) holds in the universe. Its quantifiers
say that for each \(p\in P\cap N\) there is \(q\in P\) below \(p\)
such that every \(D\in\mathcal D_P\cap N\) has \(D\cap N\)
predense below \(q\). This is exactly the required master condition.
The witness is selected separately for each \(p\); it need not be
in \(N\). No larger model's master is transferred downward, and no
generic filter or transitive model is invoked.
\end{proof}
\end{annexitem}

\begin{annexitem}{annex:p06.ppi.equivalence}{thm:p06-ppi-equivalence}{Exact quantifiers for the four equivalent tests}
\begin{proof}
Assume the all-height, all-model properness test. Choose any regular
\(\Theta>2^\kappa\) and \(A=\varnothing\). Every countable elementary
model containing \(\mathbf P,A\) is one of the models covered by
properness. This proves the fourth assertion. In fact any fixed
\(A\in H_\Theta\) could be used under this assumption.

Conversely suppose the fourth assertion holds for one such
\(\Theta,A\). The class of countably infinite elementary submodels
of \((H_\Theta,\in)\) containing \(\mathbf P,A\) is a club on
\([H_\Theta]^\omega\). For unboundedness, start with any countable
seed and add those two elements. Choose witness functions for all
formulas in the countable language of set membership, and close the
seed under these functions in countably many stages. Each stage is
countable, and the union is elementary by the Tarski--Vaught test.
For closure, the union of an increasing countable chain of these
models is countable and still contains the two parameters. A formula
with a finite tuple of parameters from that union has its tuple in
one stage; a witness exists in that stage whenever it exists in
\(H_\Theta\). The union is therefore elementary. By the fourth
assertion every model in this club has the master-production property.
This proves the third assertion without assuming a club merely from
stationarity or from the existence of isolated good models.

The trace-projection lemma proves that the third assertion implies
the second. The all-model-recognition lemma proves that the second
implies the first, at every regular \(\chi>2^\kappa\), not only at
the height where the original large models were built. This closes
the cycle. In each assertion the model test is
\((\forall p\in P\cap N)(\exists q\in P)\) with \(q\leq_P p\)
and \(q\) master for that \(N\). It is neither one condition working
below all \(p\), nor a requirement that \(q\in N\).

The large parameter can have hereditary size greater than a smaller
book-admissible \(\chi\). This causes no failure: the projected code
has hereditary size at most \(2^\kappa\), and at the smaller height
the tested model chooses a witness code of its own. The proof uses
\(\mathsf{ZFC}\), regularity and the explicit hereditary threshold;
no cardinal arithmetic assumption beyond the elementary
\((2^\kappa)^{\aleph_0}=2^\kappa\) computation is required.
\end{proof}
\end{annexitem}

\begin{annexitem}{annex:p06.ppi.cs-proper}{thm:p06-ppi-cs-proper}{Removing the full iteration datum after master production}
\begin{proof}
Fix an arbitrary stage \(\beta\leq\delta\) of the countable-support
iteration datum \(\mathcal I\) in
\cref{thm:p06-pil-extension}. The assumptions include properness of
each iterand as forced at its own stage and the actual support and
tail-name conventions of that theorem. Put
\(A=\langle\mathcal I,\beta\rangle\), and let
\(\mathbf P_\beta\) be the actual initial-segment poset.
Choose a regular \(\Theta\) large enough to contain the entire datum,
the stage and the poset, and to satisfy both the earlier lemma's
hereditary threshold for \(\mathcal I\) and the present threshold
for \(\mathbf P_\beta\). Arbitrarily large successor cardinals supply
such a regular height in \(\mathsf{ZFC}\).

For every countable \(N\prec H_\Theta\) containing
\(\mathbf P_\beta,A\), the components \(\mathcal I,\beta\) belong
to \(N\). The stage-\(\beta\) restriction of the full datum, including
its raw condition sets, quotient maps and chosen sections, is uniquely
defined from these components and also belongs to \(N\). The height
was chosen above its hereditary threshold as well. Therefore
\cref{thm:p06-pil-model-masters}, applied to this restricted datum,
supplies a master below every \(u\in P_\beta\cap N\).
This is exactly assertion four of the equivalence theorem for the
bare poset \(\mathbf P_\beta\). That theorem removes the datum and
the stage parameters and gives properness at every height admissible
for \(\mathbf P_\beta\), even when such a height does not contain
\(\mathcal I\).

Because \(\beta\leq\delta\) was arbitrary, this proves the assertion
for every initial segment and, by \(\beta=\delta\), for the full
iteration. No coordinatewise lower bound in an arbitrary proper
iterand is used here: the support-controlled proper-iteration lemma
is the independent substantive premise, and the present argument
supplies the exact parameter conversion rather than a new preservation
argument.
\end{proof}
\end{annexitem}

\begin{annexitem}{annex:p06.ppi.no-downward-inheritance}{thm:p06-ppi-no-downward-inheritance}{An atomic counterexample to inheriting a larger model's master}
\begin{proof}
Let \(\lambda\) be an uncountable cardinal and let \(P\) have a top
condition and distinct atoms \(a_\xi\), \(\xi<\lambda\). Apart from
reflexivity, the only comparisons are \(a_\xi\leq_P1_P\); different
atoms are incompatible. Every dense subset of \(P\) contains every
atom, because an atom has no proper strengthening. Work at any
book-admissible height for this poset.

The poset is proper. For a countable elementary \(K\) containing it,
an atomic \(p\in P\cap K\) is itself a \(K\)-master: it belongs to
every dense \(D\in K\), hence to \(D\cap K\), and is the sole
condition below itself. If \(p=1_P\), elementarity supplies an atom
\(a\in K\); the same argument makes \(a\) a master below \(p\).
This works for every tested model and every condition in it.

Choose a countable \(N\prec H_\chi\) containing the poset. Since
the atoms are uncountable, take an atom \(a\notin N\). A countable
elementary hull of \(N\cup\{a\}\) gives \(M\prec H_\chi\) with
\(N\subseteq M\) and \(a\in M\). Both are elementary in the same
ambient structure, so for formulas with parameters in \(N\), their
truth agrees with that in the ambient structure; hence \(N\prec M\).
As above \(a\) is an \(M\)-master.

The set \(D_{\mathrm{at}}\) of all atoms is a dense subset definable
from the poset, so \(D_{\mathrm{at}}\in N\). Its intersection with
\(N\) does not contain \(a\). Every atom in this intersection is
incompatible with \(a\), so it is not predense below \(a\).
Consequently \(a\) is not an \(N\)-master. This counterexample does
not contradict parameter invariance: another atom in \(N\) is an
\(N\)-master below the top condition. It disproves only the proposed
inheritance of an arbitrarily chosen larger-model master.
\end{proof}
\end{annexitem}

\begin{annexitem}{annex:p06.iter.proper-composition}{thm:p06-iter-proper-cited}{The general preservation statement and its all-model quantifiers}
\begin{proof}
Fix the countable-support system of Definition~\ref{def:p06-iter-recursion}
of length \(\delta\), with every stage forcing its iterand proper, and
retain its raw condition sets, bounded coordinate names, quotient maps and
chosen sections. Fix any \(\beta\leq\delta\). Choose a regular hereditary
height large enough for both this full datum and the stage poset
\(\mathbf P_\beta\), satisfying the thresholds of
Definitions~\ref{def:p06-pil-datum} and~\ref{def:p06-ppi-standing}.
For every countable elementary model containing the stage poset and the
full datum together with \(\beta\), the restricted stage datum is also
in that model. The prescribed-initial construction and
Corollary~\ref{thm:p06-pil-model-masters} therefore produce a master
strengthening of each condition in its intersection with \(P_\beta\).

Apply the fixed-parameter-to-all-model implication of
Theorem~\ref{thm:p06-ppi-equivalence}, with that full datum and stage as
the auxiliary parameter. At every height admissible for the bare stage
poset, every countable elementary \(N\) containing that poset, and every
\(p\in P_\beta\cap N\), it supplies a single \(q\leq p\) such that
\(D\cap N\) is predense below \(q\) for every dense \(D\in N\).
This is exactly Definition~\ref{def:p06-pfm-proper}; the model need not
contain the iteration datum. Thus \(P_\beta\) is proper. Since the stage
was arbitrary, every initial segment is proper, including the full
iteration at \(\beta=\delta\). The argument uses the supplied proper
fusion and parameter conversion, not the separate countable-closure
argument or inheritance of an arbitrary larger model's master.
\end{proof}
\end{annexitem}
